% 検証（F-peek）：和欧文境界グルー（CJK→LA xkanjiskip）
% peekが CJK デコーダの \fi 残骸を見ていた頃は CJK→LA 境界グルーが
% 一度も出ていなかった（pairs の ja-la 測定は abc の実幅との数値偶然で
% PASS していた）。ヘルパ CJK@X/XX/XXX/XXXX をラップして peek を
% \ifx 閉包の後に実行するよう修正し，全ルートで実トークンが見える。
\documentclass[pdflatex]{jlreq}
\begin{document}
\makeatletter
\newbox\mboxa
\def\meas#1#2#3{%
  \ifx d\jlreq@engine\global\jlreqpdf@lastclass=90\fi
  \setbox\mboxa=\hbox{#3}%
  \ifdim\wd\mboxa=#1\relax
    \typeout{PASS: #2 = \the\wd\mboxa}%
  \else
    \typeout{FAIL: #2 expected #1, got \the\wd\mboxa}%
  \fi}
% 漢字→欧文: xkanjiskip 0.25zw (pair@0@27 未定義→fallback)
\meas{17.5pt}{kanji-latin (0->27 xkanji 0.25zw)}{漢a}
\meas{27.77779pt}{kanji-abc (0->27 xkanji)}{漢abc}
% 閉じ括弧→欧文: pair@2@27 = 0.5zw (」box 5pt + 5pt + a 5pt)
\meas{15.0pt}{closebracket-latin (2->27 nat 0.5zw)}{」a}
% かな→欧文: xspc pre-only (class 15) によりグルー抑制が正の挙動
\meas{15.27798pt}{kana-latin (15->27 suppressed by xspc)}{あx}
\makeatother
% 実ページ: 漢字とabcの間に二分半の空きが見えること
\parindent=0pt
日本語とlatin textの混在。xkanjiskipが両境界に入る。
\end{document}
